% Physical unit 017; accepted Session 02.
\GordanUnitStart{017}{printed page 12}
\GordanTopFolio{12}
\[
\begin{aligned}
&=\int_{\psi_1}^{\psi_2}d\psi\,
{\begin{gathered}
\delta l'\delta l'\cos^2l(1-e^2\cos^2l)
+2\delta l\delta l'\sin l\cos l\,l'\bigl(e^2l'^2[1-e^2\cos^2l]+[1+e^2\cos^2l]\bigr)\\
+\delta l^2\bigl(e^2l'^4[\cos2l-e^2\cos^4l]+l'^2[2e^2\cos^4l-2\cos2l]\bigr)
\end{gathered}
\over \{(1-e^2\cos^2l)l'^2+\cos^2l\}^{1\frac12}}\\
&=\int_{\psi_1}^{\psi_2}d\psi\{(11)\delta l'\delta l'+2(01)\delta l\delta l'+(00)\delta l\delta l\}>0
\end{aligned}
\]
(cujus integralis' in limitibus $\delta l=0$ est).

\begin{center}\S\ 4.\end{center}

Cujus integralis elementa pro valoribus, variationi $\delta l$ vindicatis, aut positiva aut negativa fieri possunt. Totum integrale sitne positivum, an negativum, hac re definitur, utrum positivorum an negativorum elementorum summa major sit. Quoniam autem quaeritur, quatenus integrale nostrum, cujusvismodi valores integrationis sunt, positivum sit, omnis quaestio eo referenda est, quatenus integrale inveniatur, quod, cum ipsum $0$ sit, illi integrali additum, omnia ejus elementa vi positiva afficiat.

\begin{center}\S\ 5.\end{center}

Cujus generis integrale est hoc:
\[
 \int_{\psi_1}^{\psi_2}d\psi\{\delta l\delta l N\}'=
 \int_{\psi_1}^{\psi_2}d\psi\{2N\delta l\delta l'+N'\delta l\delta l\}
\]
quatenus functio $N$ ita definiri potest, ut summa
\[
 (11)\delta l'\delta l'+2[(10)+N]\delta l\delta l'+[(00)+N']\delta l\delta l
\]
intra totum spatium, $N$ nunquam infinito, integrum quadratum fiat.

Hac enim conditione
\[
 \int_{\psi_1}^{\psi_2}d\psi\{2N\delta l\delta l'+N'\delta l\delta l\}=N_2\delta l_2\delta l_2-N_1\delta l_1\delta l_1=0
\]
fit, quod variationes in limitibus evanescunt.

\begin{center}\S\ 6.\end{center}

Summa:
\[
 (11)\delta l'\delta l'+2[(10)+N]\delta l\delta l'+[(00)+N']\delta l\delta l
\]
tum existit quadratum, quum
\[
 (11)[(00)+N']=[(01)+N]^2
\]
vel integrando:
\[
 -N=(01)+(11){\alpha {\partial l'\over\partial\psi_1}+\beta {\partial l'\over\partial l_0}
 \over \alpha {\partial l\over\partial\psi}+\beta {\partial l\over\partial l_0}}
\]
qua in aequatione $\alpha$ et $\beta$ arbitrales constantes sunt.



% ---- GDL03 continuation, source PDF pages 018--027 ----
